\section{Conclusiones y Lecciones Aprendidas}

A partir de la ejecución integral de las actividades del Laboratorio N° 05, el equipo de desarrollo formula las siguientes conclusiones académicas y lecciones aprendidas:

\begin{enumerate}
    \item \textbf{Efectividad del consenso mediante Planning Poker:} La técnica demostró ser significativamente superior a la estimación empírica en horas o a la asignación arbitraria por parte de un solo integrante. La votación secreta y la revelación simultánea anularon eficazmente el \textit{efecto de anclaje} y la presión jerárquica. El debate obligatorio entre los votos extremos (el más alto y el más bajo) enriqueció el entendimiento colectivo del sistema, unificando criterios sobre la arquitectura tecnológica a utilizar en el SIGC-CUSCO (uso de \textit{Policies} de Laravel, transacciones para control de cupos y generación vectorial de certificados).
    
    \item \textbf{La dispersión de votos como indicador temprano de riesgo técnico:} Se evidenció de forma empírica que la disparidad en la escala de Fibonacci no representa un fallo en la sesión, sino un \textbf{sensor precoz de incertidumbre o ambigüedad}. Coincidentemente, las únicas tres historias donde hubo divergencias profundas (\texttt{HU-03}, \texttt{HU-07} y \texttt{HU-11}) fueron aquellas con mayor desafío de seguridad o infraestructura (código QR dinámico TOTP, aislamiento multientidad por \textit{Global Scopes} y sincronización bidireccional offline). La dispersión obligó al equipo a detenerse, dialogar y precisar el alcance antes de escribir una sola línea de código.
    
    \item \textbf{Gobernanza y trazabilidad cuantitativa en Jira Software:} La integración de las estimaciones en Jira mediante el campo \texttt{Story point estimate}, las etiquetas de clasificación (\texttt{sp-corta}, \texttt{sp-media}, \texttt{sp-compleja}) y la documentación de las rondas de votación dota al proyecto de una infraestructura de seguimiento profesional. El uso de filtros JQL permite identificar de inmediato requerimientos en riesgo y planificar la velocidad del equipo con fundamentos numéricos reales.
    
    \item \textbf{Transformación de épicas ocultas mediante el Refinamiento:} El trabajo extra evidenció el peligro de mantener historias de 13 o 21 SP en el backlog activo: concentraban más del $60\%$ del esfuerzo total del proyecto y ponían en riesgo la entrega de los sprints. La descomposición en 13 historias hijas atómicas (de 2 a 8 SP) y el aislamiento de la investigación mediante un \textbf{Spike Técnico} en \texttt{HU-11} permitió reducir la incertidumbre, logrando un Product Backlog activo de \textbf{22 historias y 96 Story Points}, donde ningún ítem supera los 8 SP, maximizando la agilidad, la verificabilidad en pruebas y la capacidad de entrega continua.
\end{enumerate}
